\documentclass[12pt]{article}
\usepackage{amsmath}
\usepackage{amssymb}
\usepackage{geometry}
\usepackage{enumitem}
\usepackage{hyperref}
\usepackage{setspace}
\usepackage{titlesec}

\geometry{a4paper, margin=1in}
\setstretch{1.5}

\begin{document}

\title{Response to Reviewers' Comments}
\date{}
\maketitle


\section*{Reviewer A}

\subsection*{General Comments}

\begin{itemize}
    \item \textit{I think that the article would benefit from some fine tuning: As it stands it is a bit vague to me in parts; I had difficulties following along some of the descriptions which I think can be fixed by thinking more about how to express the ideas and concepts in words, some restructuring and better placement of the plots (these ailments make me think this was written by a junior researcher and would benefit from an overhaul by a more experienced author).}

    \textbf{Response:} [Your response here]

    \item \textit{The figures could be placed better to be where the descriptions happen in the text. I often found myself having to turn a couple of pages to be able to follow the descriptions.}

    \textbf{Response:} [Response]

    \item \textit{I think it would be better to have one example with binary data and one with numeric data instead of two with binary data, especially since lionfish offers extra functionality for the latter.}

    \textbf{Response:} [Response]
\end{itemize}

\subsection*{Specific Comments}

\begin{itemize}
    \item \textit{Abstract: "resulting in groups who are similar but without clear separation". I think I understand what is meant, but I'm not sure this is the best way of describing it. Groups that are similar implies to me that there is no clear separation. Perhaps it is better to speak of observations that are similar and can form a cluster, but that the clusters themselves are not separated.}

    \textbf{Response:} [Response]

    \item \textit{P1: "(the cluster means)" one can use other centroids/prototypes in clustering and there are (many) clustering methods that do not use centroids/prototypes at all. Please describe whether lionfish only works for specific clusterings (e.g. k-means). Also, I'm not sure if describing the set of cluster means as "smaller number of observations" is appropriate as they are summary statistics of that larger number of observations - we still use the full set to define the objective function etc.}

    \textbf{Response:} [Response]

    \item \textit{P1: "it is usual that there are no gaps in the data". I find this is a difficult statement without backing it up; one can partition also on data with separated clusters thus, in this definition, partitioning (for gaps and non-gaps) is a superset of clustering (for gaps). Maybe it would help if you define explicitly what cluster definition you're working with - to me it sounds like you define clusters as k-means does it, namely as a Voronoi tessellation (minimum distance to any centroid).}

    \textbf{Response:} [Response]

    \item \textit{Figure 1: I'd also put the A, B, C in the text where it is described. E.g. When the correlation is high (C in Figure 1).}

    \textbf{Response:} [Response]

    \item \textit{P2: "One can see that the data is perfectly divided into four parts..." I'm not sure what to get at with this sentence and the subsequent explanation. I can use the individual features and do a k-means and get a perfect division into four clusters (k-means will almost always give a perfect partitioning if there is a unique optimum). I think you mean that the segmentation/partitioning you find in the multivariate space will not necessarily be recovered when looking at just one of the features (or any other subspace). Then, using a tour in a smart way can show properties of the multivariate space in lower dimensional subspace. I'd rewrite the paragraphs from "One can see..." to "This paper...".}

    \textbf{Response:} [Response]

    \item \textit{Figure 2: I think it would be clearer if the plots were labeled as A through C in the second row as well, as there are only the three data situations and there are no D E F. Maybe first row A (V1), ... and second row A (V2).}

    \textbf{Response:} [Response]

    \item \textit{P3: "one should be able to the separations" - is there a word missing? It also looks to me like the four bullet points would be better placed after the subsequent paragraph "While tour..."}

    \textbf{Response:} [Response]

    \item \textit{P4: "The data set must be supplied as a numeric matrix." To me this is quite is restricting. First, in R data frames are the primary data structure people use - I think the software should support this as well and not just numeric matrices. Second, marketing data are often categorical. It would make lionfish unusable in this case. Perhaps you can think of a way to incorporate categorical variables as well.}

    \textbf{Response:} [Response]

    \item \textit{P4: "The type element specifies the type of display to generate" Is there a list anywhere with all the options? I'd mention this or even better include that in the article.}

    \textbf{Response:} [Response]

    \item \textit{P4: vector of strings -> character vector in R parlance.}

    \textbf{Response:} [Response]

    \item \textit{P4: "Below this, each subset". I was wondering what subset means here. Are these the clusters? Or any selected subset of the data? Where does the subset come from?}

    \textbf{Response:} [Response]

    \item \textit{P4: This description would benefit from Figure 3 following close to the 2.1 section heading, otherwise one needs to turn the page every time to follow the description. This applies to many other plots too - it is best to have the plots on the pages where the description happens so the reader doesn't have to search for the correct figure (especially digitally).}

    \textbf{Response:} [Response]

    \item \textit{P5: "There are three additional interfaces... The first, ... the second, ...". Why not name them instead of enumerating them according to their arrangement. E.g., "The one labeled "number of bins of histograms"}

    \textbf{Response:} [Response]

    \item \textit{P5: "Helping to reduce clutter". I'm not sure what this means here.}

    \textbf{Response:} [Response]

    \item \textit{P5: "recovered by using pressing" I'm not sure what this means.}

    \textbf{Response:} [Response]

    \item \textit{P6: "These transformations are based..." To me this is an empty statement that carries no information but sounds good. If you keep it, which transformations are meant (to be straightforward to replicate one needs to know which ones)? What are "not well-established mathematical principles" and why does any of that ensure stability and efficiency?}

    \textbf{Response:} [Response]

    \item \textit{P7: "Before launching the lionfish interface, ..." This is crucial information and should be given earlier - to me this means one already needs to have a centroid based clustering result and now the subset makes sense if it is the cluster.}

    \textbf{Response:} [Response]

    \item \textit{P7: "are a summary" -> summary statistic}

    \textbf{Response:} [Response]

    \item \textit{P7: "f**f is normalized relative to each feature". I'm not sure what this means.}

    \textbf{Response:} [Response]

    \item \textit{P8: "Generally, if the features.. " How is this done?}

    \textbf{Response:} [Response]

    \item \textit{P8: "brush them". It might be good to explain what "brushing" is.}

    \textbf{Response:} [Response]

    \item \textit{Sec 3.3.: Are all these files (.png, .csv,) always generated? I think this should be controllable (I may like to save the .png but not the .csv).}

    \textbf{Response:} [Response]

    \item \textit{P9: "The responses were binarized" First, why? Second, this looks like a way of dealing with categorical and character variable inside the function (e.g., via as.numeric(as.factor()))}

    \textbf{Response:} [Response]

    \item \textit{P9: "When mostly two features.. " I don't understand this sentence at all. Please rewrite.}

    \textbf{Response:} [Response]

    \item \textit{P10: "Some of these features are more informative for the clustering."}

    \textbf{Response:} [Response]

    \item \textit{P10/11: Creating a reduced data set to look at in lionfish kind of defeats the argument you put forth in the beginning that one might overlook the clustering if looked at in reduced space. Maybe you can rewrite this.}

    \textbf{Response:} [Response]

    \item \textit{P11: "there is no objectively correct or optimal solution" -> I'm not sure what is meant here; k-means optimizes an objective function and typically has a global optimum.}

    \textbf{Response:} [Response]

    \item \textit{P12: "clusters 1 and 3 contained data points" for the reduced data set.}

    \textbf{Response:} [Response]

    \item \textit{P15: "only one representation feature" -> how was this chosen?}

    \textbf{Response:} [Response]

    \item \textit{P16: "we can assume that this group consists of tourists travelling alone as couples" is there a comma missing after alone? Because this changes the whole meaning (what about people who travel alone).}

    \textbf{Response:} [Response]

    \item \textit{P18: Sentence "While some plots, ... analyzing binary survey data..." I think the paper would benefit from at least one example illustrating the capabilities that extend far beyond the binary data type. So an example with metric data instead of two examples with binary data.}

    \textbf{Response:} [Response]
\end{itemize}

\section*{Reviewer B}

\subsection*{General Comments}

\begin{itemize}
    \item \textit{The paper does a good job of leading the reader through the steps in interactive clustering and interpretation. The lionfish software looks like a valuable contribution. My main criticisms are (i) the goals of the paper are unclear and (ii) I was unable to run the software.}

    \textbf{Response:} [Response]

    \item \textit{The paper seems to have two goals, (i) interactive visualisation for market segmentation data and (ii) description of the lionfish package. It would be helpful if this was stated clearly at the outset. The title emphasizes (i), but elsewhere in the paper (e.g., last paragraph of introduction, first sentence of Section 2, conclusion) it suggests the software development and the package lionfish is the main aim. Being clear about this would be helpful for the reader.}

    \textbf{Response:} [Response]

    \item \textit{I do not think the k-means example in the Introduction is helpful in motivating what comes later. The fact that k-means of bivariate data gives different solutions depending on correlation does not seem relevant to the content of the paper. The example does demonstrate that clusters in 2d are hidden in 1d, but I do not think a demonstration of this is necessary. I suggest instead giving an overview of market segmentation data, and the challenges and benefits of clustering such data. At the moment, the title refers to market segmentation, but we are not given an example of data of this type until Section 4.}

    \textbf{Response:} [Response]

    \item \textit{Unfortunately I cannot run the package lionfish code, interactive\_tour() produces a blank screen. Looking at the scripts for the paper, I cannot locate package pytourr.}

    \textbf{Response:} [Response]
\end{itemize}

\subsection*{Minor Issues}

\begin{itemize}
    \item \textit{Page 1, line 1: best not to refer to cluster means as observations.}

    \textbf{Response:} [Response]

    \item \textit{Page 2: "One can see that the data is perfectly divided into four parts." What does this mean? Yes, the data is partitioned, but perfectly??? It is overstating things to say there is a "distinct border between partitions" in Fig 1.}

    \textbf{Response:} [Response]

    \item \textit{Figure 2 caption is incomplete. What are colours? How does this relate to Figure 1?}

    \textbf{Response:} [Response]

    \item \textit{Since the two heatmaps in Figure 5 show different measures, I suggest using different colour schemes. Or omit B altogether, since it is not used elsewhere in the paper.}

    \textbf{Response:} [Response]

    \item \textit{Somewhere in Section 2 or 3 you need to give more explanation of the plots (aside from the heatmap, which is covered) in the lionfish GUI.}

    \textbf{Response:} [Response]

    \item \textit{Figure 5 (and elsewhere): The groupings are referred to as cluster (text) or subset (heatmap). Better to be consistent.}

    \textbf{Response:} [Response]

    \item \textit{"Projection axes with a norm < 0.1 are not shown to reduce cluttering." could be explained more clearly.}

    \textbf{Response:} [Response]

    \item \textit{When the text refers to Figures, be clear about which part of a Figure you are referring to. For example (page 13) "We can see that there is indeed overlap between clusters 1 and 3".}

    \textbf{Response:} [Response]

    \item \textit{The two examples use two different feature selection strategies, the first by inspecting a k-means of the observations and the second by clustering the features. Is there a rationale for switching between these strategies?}

    \textbf{Response:} [Response]

    \item \textit{Page 15: Feature selection: why select one feature from each cluster? They could be combined instead?}

    \textbf{Response:} [Response]

    \item \textit{Reference formatting is not consistent. Check choice of references. Ward's method not due to Murtagh and Legendre.}

    \textbf{Response:} [Response]
\end{itemize}

\end{document}